\section{Recovering inverse moduli from polyhedral data}\label{sec:pairpoly}

An embedding comes with a particularly expensive input in the repository's quantitative Jordan theorem. For a circle parametrization $f$, its inverse modulus $V$ must turn $d(s,t)\ge 2^{-k}$ into $\|f(s)-f(t)\|_\infty\ge 2^{-V(k)}$. The dimension-invariance packet also needs inverse separation, there at a specified scale. A finite table of rational affine pieces already contains this information. We extract it, including an exact best constant and an attaining pair. When the table describes a self-intersection the same calculation finds the intersection parameters.

\paragraph{Representations.}
A cell in $\mathbb Q^D$ is a bounded intersection of a finite list of closed rational affine half-spaces $h_i(z)\ge0$. Include an explicit bounding box among its rows. Equalities are pairs of rows. A pair of source points has a finite cover by such cells, on each of which two nonnegative quantities are affine:
\[
     l(z)=d(x,y),\qquad m(z)=d'(f(x),f(y)).
\]
The coordinates $z$ encode $(x,y)$ within that cell. Values agree on overlaps, and $l=0$ implies $m=0$. These are presentation conditions. In the two constructions below the coverage, value compatibility and box come directly from the input. A generic atlas used as input must supply those conditions. All finite operations below are on rational codes.

For example, on a bounded union of rational simplexes take a continuous piecewise affine map given by its vertex values, with compatible values on shared faces. Use a simplex pair with two vectors of barycentric coordinates as variables. Require each vector to be nonnegative and to sum to one. Both points and their images are affine in these coordinates. The input and output maximum norms are each the maximum of a finite list of signed affine coordinate differences. For each choice of an entry in each list, append comparisons making those entries maximal. Enumerating the choices covers all pairs, including ties. Assume a valid geometric presentation, so two barycentric descriptions of one source point give the same image. The same construction works for a domain consisting of one simplex and for triangulated cubes.

For a polygonal circle let $0=b_0<b_1<\cdots<b_N=1$ be rational breakpoints. Give $P_i\in\mathbb Q^m$ for $0\le i\le N$ with $P_N=P_0$, $m\ge1$. On $[b_i,b_{i+1}]$ use affine interpolation. Symmetry allows $s\le t$. In each rectangle of parameter intervals impose that order and split it by $t-s=1/2$. The two formulas for $l$ are $t-s$ and $1-t+s$. For each possible signed output coordinate impose that it dominate all the others and use it for $m$. Thus $D=2$ independent of the number of edges and target dimension. The pair $(0,1)$ denotes one circle point and gives $l=m=0$.

\begin{lemma}[Finite vertex interpolation]\label{lem:pairvertex}
For a cell $H$, enumerate all choices of $D$ rows, solve their equality systems with nonzero determinant, and retain solutions satisfying every row. Every rational $z\in H$ has a primitive recursively constructed binary tree of convex combinations of retained points, of depth at most $D$. Duplicate retained points do no harm.
\end{lemma}
\begin{proof}
At $z$ mark a row active when $h_i(z)=0$. Compute the rank $r$ of the active linear parts by scanning nonzero minors (the empty minor is one). If $r=D$, the independent active rows yield $z$ by Cramer's rule, so this point is on the list. Otherwise choose a nonzero rational vector $w$ annihilated by all active rows. Explicitly choose an $r$-by-$r$ invertible row-column submatrix, set one free coordinate to one and the other free coordinates to zero, and solve its $r$ equations for the pivot coordinates. Every remaining active row is in the span of the pivot rows: after their elimination any nonzero free entry would give a nonzero $(r+1)$-minor. This also covers $r=0$.

Along $z+tw$ an inactive row has positive slack $h_i(z)$ and directional slope $a_i\cdot w$, where $a_i$ is its linear part. Among negative slopes take the least positive $h_i(z)/(-a_i\cdot w)$ to get $t_+$. Along $-w$ obtain $t_-$ by the same rule. A nonzero coordinate of $w$ and its two box rows ensure both lists are inhabited. An active box face would force that coordinate of $w$ to vanish. The two endpoints $z^+=z+t_+w$ and $z^-=z-t_-w$ satisfy every row, retain all old active rows and activate a new row with nonzero slope on $w$. Such a row is independent of the old rows. Their convex combination with respective weights $t_-/(t_-+t_+)$ and $t_+/(t_-+t_+)$ is $z$. In a depth-$D$ tree split unmarked nodes in this way, padding finished leaves by their own copies if desired. Along every branch the active rank increases at each genuine split, so the leaves belong to the enumerated list. Substitution proves the convex-combination identity at every node. All minima used here are over explicit row lists.
\end{proof}

\begin{theorem}[Sharp distortion or a rational collision]\label{thm:pairratio}
Suppose at least one pair has positive source distance. On all retained points in all pair cells with $l(v)>0$, compute
\[
       c=\min_v\frac{m(v)}{l(v)},\qquad L=\max_v\frac{m(v)}{l(v)}.
\]
Then $c\,d(x,y)\le d'(f(x),f(y))\le L\,d(x,y)$ for every rational pair, and both bounds are attained at rational pairs of positive source distance. If $c=0$ the attaining pair is a collision. If $c>0$ it gives quantitative injectivity. If there is no positive-distance pair, the scan is empty and every distance is zero.
\end{theorem}
\begin{proof}
A positive $l(z)$ is a convex combination of nonnegative leaf values, hence some leaf has positive $l$. The scan is inhabited. The inequalities hold at each retained point of positive $l$ by the two comparisons in the scan. At every retained point of zero $l$, the consistency condition gives $m=0$. Thus no division on the diagonal is required. On each cell both inequalities are affine inequalities. Propagate them up the tree by multiplying each inequality by its nonnegative weight and adding. The result at the root holds for the specified pair. As every retained point is feasible, the two winning pairs are genuine and realize the two ratios. A positive-distance collision forces $c=0$ by the lower bound. For named pairs the same inequalities follow at each requested precision by rational approximation, since both metrics and the specified continuous PL functions respect limits. In the rational PL constructions the upper bound provides an explicit continuity rate for this last passage.
\end{proof}

\paragraph{Discharging the repository hypotheses.}
Write the positive $c=a/b$ and $L=p/q$ with positive denominators. Take $j_-$ least with $2^{j_-}a\ge b$, scanning through $b$, and $j_+$ least with $2^{j_+}q\ge p$, scanning through $p$. The inequalities $2^b\ge b$ and $2^p\ge p$ justify the budgets. For the Jordan packet take
\[
      V(k)=k+\max(8,j_-),\qquad U(k)=k+\max(8,j_+).
\]
The rational evaluator is simply affine interpolation on the first interval containing the parameter. The inequalities of the theorem imply both modulus premises in the packet at rational parameters, even before its allowance for two approximation errors. For a PL cube the dimension-invariance packet asks that an image discrepancy (including evaluation error) below $2^{-a_0}$ imply source distance at most $1/2$. Take $a_0=1+j_-$. An actual image discrepancy below $2^{-a_0}\le c/2$ forces $d(x,y)<1/2$ by the lower bound. The extra nonnegative evaluation-error term in that packet preserves the implication for the exact rational evaluator. A crossing circle returns explicit unequal parameters instead. This supplies numerical regularity data to the existing mathematical proofs. Their remaining parity and incidence derivations in Lean are separate obligations.

\paragraph{Two examples and a perturbation.}
Traverse the unit square in four equal parameter intervals. Its exact $(c,L)$ is $(2,4)$. On one edge the ratio is $4$; on adjacent edges let the two distances to the shared vertex be $a,b$ in edge-length units. The shorter parameter distance is at most $(a+b)/4$, and the image distance is $\max(a,b)$. On opposite edges the image distance is at least one and the parameter distance is at most $1/2$. The opposite vertices attain ratio two. Traversing the bow-tie $(0,0),(1,1),(0,1),(1,0),(0,0)$ in the same intervals gives $c=0$ at $s=1/8$, $t=5/8$. If a map $g=P+R$ satisfies $\|R(s)-R(t)\|_\infty\le e\, d(s,t)$, the triangle inequality gives lower and upper constants $c-e$ and $L+e$ for $g$. For a curved example use the square and $R(s)=(s(1-s),0)$. For $s\le t$ the factorization of the difference gives $|R_1(t)-R_1(s)|\le t-s$. Applying the same bound on the two endpoint intervals and $R(0)=R(1)$ also gives $1-t+s$. Thus $e=1$ and the curved circle has bounds $(1,5)$, with fully rational evaluation at rational parameters.

\paragraph{A version for uniform approximation.}
Suppose a polygon $P$ is uniformly within $e$ of a curve at the current demand $k$. Intersect its pair cells with the single new half-space $l\ge 2^{-k}$. Minimize the affine image distance over the retained vertices, obtaining $h$. If these cells are empty, the separation condition at this scale is vacuous by Lemma~\ref{lem:pairvertex}. Otherwise every eligible image pair of the curve is separated by at least $h-2e$, by the triangle inequality. A check $h>2e$ now produces the inverse exponent by the same finite denominator bound as above. This test can succeed even if $P$ crosses itself at a smaller parameter separation. A PR choice of such table and error for each $k$, with checks passing, suffices for a uniform inverse modulus (take prefix maxima and pad). A failed test establishes only that this particular estimate is insufficient.

\begin{proposition}[Why a forward modulus alone gives no uniform PR extraction]\label{prop:delayedcurves}
There is a single PR evaluation scheme for a family of embedded circles with a common forward Lipschitz bound for which no PR function of the family index supplies even $V(1)$ uniformly. This is an external assertion about the standard natural numbers.
\end{proposition}
\begin{proof}
Enumerate all syntactic unary PR definitions, treating invalid codes as the zero function. An ordinary machine $M(e)$ evaluates definition $e$ on $e$ and then spends more than the result in additional steps. For each standard $e$ the computation terminates: evaluate compositions by their constituents and recursion by a finite loop. Let its running time be $T(e)$. For any unary PR function $g$ take an index $e$ of $g$; then $T(e)>g(e)$.

Let $P$ be the unit-square circle above and set $f_e(s)=2^{-T(e)}P(s)$. At demand $n$ run exactly $n$ steps of the machine simulation. If it halted, its time is available and the evaluator returns $2^{-T(e)}P(s)$. Otherwise return the zero vector. The bounded simulation is PR. The uniform error is at most $2^{-n}$, since $\|P\|_\infty\le1$. Each curve is a positive multiple of an embedding and all have forward Lipschitz bound 4. Opposite vertices at parameters $0$ and $1/2$ have distance exactly $2^{-T(e)}$. An inverse exponent $g(e)$ at $k=1$ therefore satisfies $g(e)\ge T(e)$ for every $e$, contradicting the preceding diagonal argument. The example concerns uniformity over presentations; each fixed index has a valid numerical modulus. It makes no claim that iPRA proves termination of $M$ uniformly. The vertex algorithm instead accepts the actual rational PL table, whose denominators contain the missing quantitative information.
\end{proof}

\paragraph{Arithmetic and Lean obligations.}
The PR loops enumerate row subsets (bit masks below $2^H$ for $H$ rows), column subsets and permutations for determinants. Cramer's rule uses rational products and bounded comparison, with zero divisors tested first. Rank uses a finite maximum. For the interpolation proof use the complete binary tree of depth $D$, with at most $2^{D+1}-1$ positions. At level $a$ the invariant is the conjunction over its node indices that each point is feasible and its rank is at least $\min(a,D)$; finished vertices can be repeated. Additional conjuncts check each parent interpolation identity and the nonnegative rational weights. At a node the row tests, minor scans, determinant expansion and endpoint computations are finite identities or comparisons; determinant row elimination follows by multilinearity and alternating cancellation of the bounded permutation sum. The rank-growth check can also be proved by adjoining the new row to the pivot submatrix: eliminating the old pivot columns of that row leaves a nonzero free entry, since it has nonzero dot product with $w$. Backwards level induction on the bounded conjunction of the two affine inequalities propagates the leaf bounds. Node states are PR functions of the input and level, so even enormous intermediate numerators do not introduce an unbounded existential in the invariant. The chart and queried rational pair stay fixed during these inductions. The resulting universal closure over numerical codes and the two alternative numerical outputs are legitimate iPRA targets. There is no induction on a statement quantifying over all convex sets or real vectors.

The executable \texttt{checks/polygon\_modulus.py} implements the two-variable algorithm. It checks both extrema by independent interpolation for eight examples, including crossing, retraced and constant maps. In \texttt{lean/PairBounds.lean}, \texttt{Mix} is the inductively defined list-vertex/convex-split relation on rational pairs $(l,m)$. The checked lemma \texttt{lower\_of\_mix} propagates a finite Boolean lower-bound test through a supplied tree. The upper inequality follows by carrying $L l-m\ge0$ through the same affine-combination rule. The file also checks sharpness at an attained pair and the two perturbation inequalities. The actual object-language file \texttt{lean/PairBoundsIPRA.lean} proves \texttt{closed\_mix\_lower}. Its matrix is
\[
 Au\le BU\ \longrightarrow\ Av\le BV\ \longrightarrow\
 A(ru+sv)\le B(rU+sV)
\]
with all eight number variables universally closed. Clearing denominators in one convex split gives exactly this implication, because distances and weights are nonnegative. For the upper bound exchange the two distances and the two slope factors. To finish the full \texttt{Prf PRAAx} endpoint, implement the row, determinant and tree computations as fixed PR symbols, prove the rank/active-row invariant above, apply this split rule level by level and reflect the bounded argmin and argmax checks. The closed split rule is a checked object derivation. The rational host lemmas, constructive interpolation proof and numerical runs have the individual scopes stated here.
